\chapter{Consonance as a Relational Function}

\epigraph{\emph{Consonance does not belong to the isolated ratio alone.}}{}


If the previous chapter argued that harmonic information is best understood in relational terms, then consonance can no longer be treated as a self-sufficient essence lodged inside a ratio taken in abstraction from every medium, body, and context. The ratio remains indispensable. It would be pointless to deny that simple proportional relations recur across acoustics, dynamical systems, and wave-bearing systems of many kinds with unusual explanatory force. Yet that force is not adequately described by saying that consonance simply inheres in numerical simplicity itself. A ratio does not sound, radiate, interfere, stabilize, or become salient in a vacuum. It does so only within a concrete field composed of modal constitution, physical transmission, coupling regime, receptive organization, and contextual framing. For that reason, the present chapter proposes a stricter formulation: consonance should be understood as a relational function rather than as an intrinsic property of the interval considered in isolation.

This shift is not a denial of structure but a correction of emphasis. HIT does not dissolve consonance into pure relativism, nor does it reduce it to arbitrary taste. It instead claims that the ratio names a privileged relational constraint whose consequences depend on the system in which it appears. One may summarize the point provisionally as \emph{Consonance = F(ratio, modal or spectral constitution, medium, coupling regime, receiver, context).} In the acoustic case, that modal constitution appears as timbre and overtone profile; elsewhere, it may appear as mode structure, propagation condition, or resonant topology. Such a formula should not be read as a finished equation, but as a disciplinary warning. It reminds us that no account of consonance is conceptually complete if it abstracts away from the constitution of the interacting modes, the resonant properties of the medium, the dynamics of the receiving system, and the practical or perceptual horizon in which the relation is encountered.

The term \emph{consonance} is retained here because the acoustic and musical traditions gave this family of problems one of its richest and most durable vocabularies. But the concept, as HIT develops it, should be read more broadly. It names a kind of resonant fit, mode compatibility, or low-conflict coordination that becomes especially legible in the sonic domain without being exhausted by it. The task of this chapter is therefore double: to preserve what the sound-focused traditions discovered, and to prevent those discoveries from being mistaken for the limits of the field itself.


\section{Against the isolated essence of the ratio}


The strongest traditional intuition about consonance is easy to state: simple ratios tend to generate more stable, more coherent, or more easily integrated organizations than complex ones. In acoustics this often appears first as a matter of sounding intervals, but the intuition is not reducible to the auditory case alone. It is grounded in a long history of observation about periodic coincidence, spectral overlap, recurrence, and perceptual organization. From classical acoustics through later psychoacoustics, relations such as $2{:}1$, $3{:}2$, and $4{:}3$ have repeatedly appeared as privileged structures rather than as arbitrary mathematical curiosities \parencite{helmholtz_1863,plomp_1965}. Even when the explanatory models differ, the recurrence of these ratios is too persistent to be dismissed. The problem lies elsewhere. The problem begins when one mistakes this recurrent privilege for a complete ontology of consonance.

The insufficiency of that move becomes clear as soon as constitution enters the picture. Sethares showed with unusual clarity, in the sonic case, that consonance is not fixed by interval alone, but by the relation between intervallic structure and the spectral organization of the sounding bodies involved \parencite{sethares_2005}. The same numerical ratio can behave differently under different timbral conditions because what matters is not merely proportional abstraction, but the alignment or conflict of partials, beating structure, and spectral fit. The broader lesson extends beyond acoustics: a relation acquires its effective character only through the modes that realize it and the medium in which those modes interact. This does not abolish the importance of the ratio. It relocates it. The ratio becomes a structural organizer whose perceptual and dynamical consequences depend on the ecology in which it is instantiated. In that sense, the ratio remains necessary but ceases to be sufficient.

Classical roughness theory already pointed in this direction. Plomp and Levelt's critical-band model explained much of consonance and dissonance not through number mysticism but through interference among partials falling within or outside sensitive auditory bandwidths \parencite{plomp_1965}. Helmholtz likewise rooted consonance in the physiology of hearing and in the interaction of overtones rather than in arithmetic purity alone \parencite{helmholtz_1863}. These traditions are sometimes read as dismantling the older ratio-centered view. A better reading is that they complicate it. They show that consonance emerges where proportional relation, spectral constitution, and reception intersect. If one keeps only the ratio, the phenomenon becomes abstract to the point of distortion. If one keeps only local physiology, one loses the structural persistence that made the phenomenon visible in the first place. The acoustic archive is therefore best read as a particularly articulate instance of a more general problem of relational fit.

The deeper conceptual issue is that an isolated ratio easily becomes a metaphysics of presence. Once detached from medium and receiver, it begins to appear as though its significance were fully present in itself, prior to any field of manifestation. This is precisely the kind of simplification that the broader philosophical background of the project resists. Heidegger's emphasis on disclosed worldhood suggests that significance appears within a horizon rather than as an intrinsic content of detached objects \parencite{heidegger_1927}. Derrida's critique of self-presence presses the same problem from another side: what seems self-identical often depends on a field of difference and articulation that it cannot contain within itself \parencite{derrida_1967}. Nietzsche's suspicion toward reified concepts adds a final caution: once a category hardens into a supposedly transparent truth, its historical and interpretive scaffolding tends to disappear from view \parencite{nietzsche_1873}. In the present context, these lines of thought converge on a simple warning: one should not mistake the privileged status of harmonic ratios for the claim that consonance is exhausted by numerical simplicity alone.

For HIT, then, the ratio must be retained without being absolutized. It is best understood as an attractor of structured coordination, not as an autonomous essence. The question is no longer whether certain ratios matter. They plainly do. The question is how they matter, through what material and perceptual mediations, and under what conditions they become more or less consonant for a given system. Once posed in that way, consonance ceases to look like a timeless property of intervals and begins to look like a dynamic outcome of relation.


\section{Consonance as dynamic stability rather than rigid perfection}


The relational reformulation of consonance becomes sharper when one replaces the image of mathematical purity with that of dynamic stability. In many domains, harmonic simplicity matters not because it imposes rigid identity, but because it opens wider basins of coordination. Kuramoto-style models make this especially clear: simple rational relations tend to support more robust phase-locking regimes and wider Arnold tongues than more complex relations do \parencite{kuramoto_1984,pikovsky_2001,strogatz_2000}. Recent work on the forced Kuramoto model with matrix coupling strengthens this point by showing that the architecture of coupling itself can open multiple resonant regimes beyond the standard $1{:}1$ case, yielding Arnold tongues and devil's staircases whose structure depends on the coupling matrix rather than on scalar forcing alone \parencite{costa_deaguiar_2026}. What is privileged in such systems is not immobility, but coordinated persistence under perturbation. A stable ratio is therefore not best imagined as a frozen number. It is better understood as a dynamically favored regime toward which interacting oscillators can be drawn and within which they can remain coherent without becoming identical.

This distinction matters because a great deal of intuitive discourse around harmony still carries the residue of a perfectionist picture. On that picture, the closer a relation approaches an ideal numerical purity, the more consonant it must be, as if the goal of harmonic organization were the elimination of deviation altogether. But dynamical systems rarely reward that kind of rigidity. Biological, acoustic, and ecological coordination often depend on a margin of tolerance, flexibility, and adaptive response. Partial synchronization, metastability, and robust recurrence are frequently more informative than exact repetition. Acebron and colleagues made this plain in their review of Kuramoto-class systems, where partial coherence and clustered coordination appear not as failures of order but as characteristic forms of it \parencite{acebron_2005}. Glass and Mackey showed something similar in physiological systems, where rhythmic order is often preserved through flexible locking rather than strict identity \parencite{glass_1988}. Driven media and condensates show a comparable lesson in another idiom: periodic forcing does not yield order by abolishing fluctuation, but by producing stable patterning regimes within it \parencite{engels_2007,nguyen_2019}.

Maturana and Varela's concept of structural coupling gives this dynamic picture a biological generalization. Structural coupling names a history of recurrent interactions through which system and environment become mutually congruent without ever becoming identical or exchanging information in the naive transmissive sense \parencite{maturana_varela_1980}. Arnold tongues can be read as one formal image of that problem for oscillatory systems: some relations offer wider and more stable basins of capture than others, which means that certain couplings are easier to enter and sustain. In that restricted sense, consonance can be reformulated as successful coupling rather than as static intervallic purity. A consonant relation is one that recurrently supports coordinated persistence under the material conditions of the system in question.

This reformulation also clarifies what earlier sections called stability without rigidity. Structural coupling does not require perfect sameness; it requires a recurrent compatibility between the organization of the system and the perturbations it encounters. That makes room for drift, tolerance, and adaptive variation without forfeiting the ratio-like regularities that make coordination possible in the first place. Di Paolo's later account of precarious autonomy helps sharpen the point further: when an organized system is fragile, the regimes that better sustain congruent interaction matter normatively, not just descriptively \parencite{dipaolo_2005}. The analogy should not be overdriven. Kuramoto oscillators are not organisms, and biological coupling includes memory, plasticity, and history in ways generic mode-locking models do not. But the formal convergence is strong enough to justify a conceptual shift: consonance is best treated as a relational achievement of successful coupling rather than as an intrinsic jewel hidden inside isolated numerical ratios.

From this perspective, the ratio should be treated as an attractor rather than a destiny. A consonant regime is not one in which all variation has been suppressed, but one in which the relevant relations can maintain coherent organization across time despite local fluctuations. That is why recurrence-based approaches are so suggestive here. Trulla and colleagues frame consonance not only as a matter of cultural naming or auditory preference, but as a property of time-domain organization: simpler relations tend to produce stronger and more legible recurrence structure \parencite{trulla_2018}. Recurrence does not reward absolute stasis. It rewards patterned return. This is exactly the kind of criterion that allows HIT to speak of natural harmony without collapsing into geometrical fetishism.

One can therefore formulate the chapter's central correction as follows: consonance is not the same thing as exactness. A relation may be consonant because it is sufficiently stable to sustain coordinated organization and sufficiently flexible to absorb perturbation without disintegrating. That is a stronger and more plausible thesis than the idea that consonance belongs to the perfectly rigid realization of an abstract ratio. It also helps explain why living and material systems can exhibit harmonic organization without ever instantiating ideal mathematical purity. The relevant question is not whether a system has achieved a pristine ratio in the abstract. The question is whether the relation organizes interaction in a way that remains coherent, recurrent, and low-conflict within the concrete dynamics of the system.

This also clarifies the role of deviation. Deviation from an ideal ratio is not automatically evidence against natural harmony. Under many conditions, a bounded deviation may express the adaptive capacity of the system rather than its failure. To say that harmony is \enquote{sufficiently stable, sufficiently dynamic} is therefore not a rhetorical compromise. It is an ontological claim about the kind of order consonance names. The aim is not perfect rigidity, which in many natural contexts would amount to fragility or even functional death. The aim is a regime of relation within which stability and responsiveness are co-sustained. Once framed this way, consonance begins to look less like an arithmetic relic and more like a dynamical achievement.


\section{Natural harmony and tempered harmony}


The distinction between natural harmony and tempered harmony can now be restated more carefully. Here the musical domain remains central, because it is one of the historical laboratories in which these problems were named with exceptional density. The first term names a family of relations that arise from the resonant organization of physical systems, especially from overtone structure and the patterned constraints of coupled oscillation. The second names a historically elaborated cultural technology that redistributes intervallic space in order to secure transpositional flexibility, modulatory reach, and instrumental compatibility. Both belong to the history of music and sound. They should not be confused, and they should not be turned into caricatures of one another.

The harmonic series is the most familiar index of natural harmonic organization. When a resonant body vibrates, it does not generally produce only a fundamental frequency; it also produces a hierarchy of partials whose proportional relations are far from arbitrary. Those relations make available intervallic neighborhoods that recur across instruments, voices, and material media. Partch's \emph{Genesis of a Music}, first published in 1949 and reissued in its second edition in 1974, built much of its critique of equal temperament on the conviction that these relations were not merely historical curiosities but traces of a deeper acoustic order that Western musical standardization had partly obscured \parencite{partch_1949_1974}. Tenney's historical work likewise makes clear that the concepts of consonance and dissonance have always moved within a tension between acoustic inheritance and cultural systematization rather than belonging wholly to one side or the other \parencite{tenney_1988}. HIT enters that tension neither to abolish culture nor to canonize raw acoustics, but to keep the distinction conceptually visible.

Temperament, and especially twelve-tone equal temperament, should be understood in this light as an ingenious compromise rather than as a falsification of nature. It regularizes intervallic space by distributing deviation across the octave so that modulation and key mobility become practicable within a standardized grid. Its power is historical, technical, and expressive. But that power is purchased through a transformation of the underlying relation between number and resonance. The equal-tempered system does not preserve all natural proportional relations exactly; it replaces them with a more uniform architecture optimized for other purposes. What one gains is portability across tonal centers. What one loses is strict fidelity to the local resonant logic of particular simple ratios. As Nicol\'as Ech\'aniz, president of Asociaci\'on Civil AlterMundi and a collaborator in this program, has observed in conversation, much of the music we make remains slightly out of tune with how the universe sounds. Taken carefully, the phrase should not be heard as cosmological grandiosity. It names the gap between a physical regime of resonance and a culturally useful regime of intervallic standardization.

Seen from this angle, modulation itself acquires two distinct meanings. In a natural-harmonic register, modulation may be understood as a change of center, perspective, or resonant anchoring within an already structured field of relations. In the tempered register, by contrast, modulation is made possible by the relative uniformity of the grid; one moves across key centers because the system has deliberately flattened local differences in interval quality. Neither regime is illegitimate. They simply solve different problems. The first remains closer to the uneven texture of resonance. The second privileges combinatorial freedom and transportability. HIT needs this distinction because it cannot claim that every musically successful organization is an expression of natural harmonic structure. Nor can it pretend that the physical organization of resonant systems is reducible to a cultural lattice designed for convenience and expressive expansion.

This distinction also helps prevent a familiar mistake. Once natural harmony is opposed too simplistically to tempered harmony, the former starts to look like authenticity and the latter like corruption. The stronger claim is more specific: natural harmony names a set of physically grounded relational tendencies; tempered harmony names a historically constructed strategy for redistributing and managing those tendencies. The question for HIT is not which one is morally superior. The question is what kinds of informational, perceptual, and dynamical consequences follow from each mode of organization.


\section{Each physical system has its own natural harmony}


If consonance is relational and dynamically situated, then natural harmony cannot be treated as a single universal template imposed identically across all systems. The very phrase must be localized. A resonant cavity, a photonic medium, a Bose-Einstein condensate, a wooden instrument, a vocal tract, a thoracic cavity, and a forest soundscape do not expose precisely the same affordances even when some proportional regularities recur across them. Their material constraints differ; their coupling possibilities differ; the receivers that interact with them differ as well. What counts as stable, legible, or low-conflict relation in one system may not transfer cleanly to another without qualification.

This localized view is entirely compatible with the idea that simple harmonic relations often remain privileged. What changes is the meaning of \enquote{privileged.} It no longer means universally sovereign in the same form everywhere. It means that within a given physical ecology, certain relations become recurrently favored because they fit the resonant and organizational properties of that ecology. Sethares's timbre-sensitive account already points in this direction in acoustics, since the consonance of an interval depends on the spectra that realize it \parencite{sethares_2005}. The same proportional relation may become more or less consonant as the material and modal constitution of the system changes. To speak of a system's natural harmony is therefore to speak of the relation between its structural tendencies and the concrete way those tendencies are materially instantiated.

This formulation protects HIT from a naive universalism. It allows one to say that natural harmonic organization is real without claiming that one fixed scale or one fixed set of ratios governs every meaningful case in an identical way. It also preserves the relevance of embodiment. A receiving system does not encounter harmonic relations from nowhere. It encounters them through a body, a medium, and a history of coupling. In that sense, natural harmony is always local twice over: local to the sounding system and local to the receiving system. The consonance of a relation emerges where these localities can meet in a sufficiently coherent way.

The consequence reaches beyond this chapter. Once natural harmony is understood as situated, the task of HIT becomes to identify, describe, and compare the resonant regimes through which concrete systems achieve relative coherence. Whether the system is an instrument, a body, an ecosystem, or an artificial device, harmony must be sought in the relations that its material organization can actually sustain. This orientation gives unity to the chapters that follow: they ask how different systems organize, register, and respond to patterns of resonant order, and how those patterns might be measured, modeled, and eventually designed with greater precision.


\section{A minimal relational proposition about consonance}


The chapter's argument may now be condensed into a minimal proposition. Consonance is not the intrinsic property of a ratio considered apart from its realization. It is the capacity of a relation among oscillatory, wave-bearing, or mode-organized processes to sustain a coherent, low-conflict, and differentially interpretable organization within a concrete system. The ratio remains central because it names a structural constraint with unusual power to organize recurrence and coordination. But that power is actualized only through constitution, medium, coupling regime, receptor, and context. Consonance, in other words, belongs neither to number alone nor to arbitrary judgment alone. It belongs to the relational achievement of stability within a situated field.

This proposition should be read with the same distinctions introduced at the end of the previous chapter. The observation is that simple harmonic relations repeatedly appear as privileged sites of recurrence, locking, and perceptual coherence. The hypothesis is that consonance is best understood as a relational function whose value depends on the organization of the system in which the relation appears. The inference is that HIT must reject both naive Pythagorean essentialism and flat cultural conventionalism. The former mistakes ratios for complete causes. The latter risks obscuring the non-arbitrary regularities that keep returning across physical and biological domains.

Once framed this way, the next step becomes unavoidable. If consonance is a relational function rather than a self-sufficient property, then the central claims of HIT can no longer remain implicit. They must be stated explicitly, ranked by evidential status, and exposed to possible failure. That is the task of the next chapter, which turns from conceptual clarification to the formal enumeration of the program's central hypotheses.
